\begin{table}[p]
\centering
\caption{\\ Gold clause exposure and investment by firm size, cash, and book leverage } \scriptsize \renewcommand{\arraystretch}{1.15}

\label{tabapp:constraint}
{
    \def\sym#1{\ifmmode^{#1}\else\(^{#1}\)\fi}
    \begin{tabular}{l*{3}{D{.}{.}{-1}}}
        \toprule
                            & \multicolumn{1}{c}{Small} & \multicolumn{1}{c}{Low cash} & \multicolumn{1}{c}{High leverage} \\
                            & \multicolumn{1}{c}{(1)} & \multicolumn{1}{c}{(2)} & \multicolumn{1}{c}{(3)} \\
        \midrule
        Q & 0.095\sym{***} & 0.096\sym{***} & 0.096\sym{***} \\
         & (0.017) & (0.017) & (0.017) \\
        \ensuremath{\tilde{d}} & 0.012 & -0.068\sym{**} & -0.048\sym{**} \\
         & (0.024) & (0.023) & (0.019) \\
        \ensuremath{\tilde{d}} \ensuremath{\times \text{ I}} & -0.103\sym{***} & 0.042 & 0.010 \\
         & (0.029) & (0.038) & (0.058) \\
        1933 \ensuremath{\times \text{ I}} & 0.010 & -0.053\sym{**} & -0.034 \\
         & (0.018) & (0.019) & (0.032) \\
        1934 \ensuremath{\times \text{ I}} & -0.000 & -0.023 & -0.033 \\
         & (0.015) & (0.016) & (0.028) \\
        After \ensuremath{\times \text{ I}} & -0.009 & -0.022 & -0.009 \\
         & (0.015) & (0.015) & (0.025) \\
        1933 \ensuremath{\times \tilde{d}} & -0.036 & -0.086\sym{***} & -0.061\sym{***} \\
         & (0.022) & (0.022) & (0.016) \\
        1934 \ensuremath{\times \tilde{d}} & -0.016 & -0.053\sym{**} & -0.058\sym{***} \\
         & (0.021) & (0.019) & (0.013) \\
        After \ensuremath{\times \tilde{d}} & -0.034 & -0.013 & 0.000 \\
         & (0.023) & (0.016) & (0.016) \\
        1933 \ensuremath{\times \tilde{d} \times \text{I}} & -0.045 & 0.068\sym{*} & 0.012 \\
         & (0.039) & (0.033) & (0.057) \\
        1934 \ensuremath{\times \tilde{d} \times \text{I}} & -0.053\sym{*} & 0.027 & 0.060 \\
         & (0.028) & (0.030) & (0.048) \\
        After \ensuremath{\times \tilde{d} \times \text{I}} & 0.045 & 0.023 & -0.017 \\
         & (0.027) & (0.031) & (0.044) \\
        \midrule
        Firm FE & \multicolumn{1}{c}{Yes} & \multicolumn{1}{c}{Yes} & \multicolumn{1}{c}{Yes} \\
        Year FE & \multicolumn{1}{c}{Yes} & \multicolumn{1}{c}{Yes} & \multicolumn{1}{c}{Yes} \\
        \ensuremath{R^2} & 0.240 & 0.239 & 0.239 \\
        Observations & \multicolumn{1}{r}{7,066$\phantom{000}$} & \multicolumn{1}{r}{7,066$\phantom{000}$} & \multicolumn{1}{r}{7,066$\phantom{000}$} \\
        \bottomrule
    \end{tabular}
}

\vspace*{3mm} \justifying \noindent
\scriptsize{\textit{Notes.} This table tests whether debt overhang effects vary by firm characteristics. It reports results from panel regressions of net investment on $Q$, gold clause exposure $\tilde{d}$, period $\times$ $\tilde{d}$ interactions, and triple interactions with a firm characteristic indicator $I$, where the pre-abrogation period (1926--1932) is the omitted category. In column 1, $I$ equals one if a firm has below-median asset size among $\tilde{d} > 0$ firms in 1930. In column 2, $I$ equals one if a firm has below-median cash/assets among $\tilde{d} > 0$ firms in 1930. In column 3, $I$ equals one if a firm has above-median book leverage among $\tilde{d} > 0$ firms in 1930. All regressions include firm and year fixed effects. All variables are winsorized at the 0.5\% and 99.5\% levels within each year. Standard errors in parentheses are two-way clustered by firm and year. $^{*}p<0.10$, $^{**}p<0.05$, $^{***}p<0.01$.}
\end{table}